% ============================================================================
% GÉNÉRÉ par scripts/exemples-guide.ts — NE PAS ÉDITER À LA MAIN.
% Régénérer : npm run exemples (chaque chiffre est vérifié contre le moteur).
% ============================================================================
\begin{table}[ht]
\centering
\small
\begin{tabular}{r r r}
\toprule
Revenu de travail (\$) & Barème seul & TEMI réel \\
\midrule
\num{0} & \pc{0} & \pc{18.3} \\
\num{5000} & \pc{0} & \pc{35.4} \\
\num{10000} & \pc{0} & \pc{59.3} \\
\textbf{\num{15000}} & \textbf{\pc{0}} & \textbf{\pc{82.7}} \\
\num{20000} & \pc{26.1} & \pc{44.9} \\
\num{25000} & \pc{26.1} & \pc{51.1} \\
\num{30000} & \pc{26.1} & \pc{33.5} \\
\num{35000} & \pc{26.1} & \pc{43.4} \\
\textbf{\num{40000}} & \textbf{\pc{26.1}} & \textbf{\pc{80.6}} \\
\num{45000} & \pc{26.1} & \pc{54.8} \\
\num{50000} & \pc{26.1} & \pc{57.6} \\
\num{55000} & \pc{26.1} & \pc{60} \\
\textbf{\num{60000}} & \textbf{\pc{36.1}} & \textbf{\pc{64.9}} \\
\bottomrule
\end{tabular}
\caption{TEMI 2025 du ménage type M6 (famille monoparentale, un enfant de
3~ans en garde subventionnée), par tranche de \D{5000} : chaque ligne donne
le taux effectif sur la tranche de revenu qui débute au montant indiqué. En
gras, les tranches en \emph{trappe} (TEMI $>$ \pc{60}). La colonne
\og barème seul \fg{} est le taux marginal des seules tables d'imposition
au revenu imposable correspondant (zéro sous le montant personnel de base).}
\label{tab:temi-m6}
\end{table}
